\section{Reed--Solomon lattices and ideal cosets}
\label{sec:descent}

The first step of the reduction stores an X3C selection vector
$\boldsymbol\xi\in\{0,1\}^n$ in designated coordinates of a longer vector
$\mathbf x\in\{0,1\}^{\F_q}$.  Evaluation at $\zeta$ maps $\mathbf x$ to
$[\mathbf x]=\sum_{a\in\F_q}x_a\zeta^a\in\mathcal{O}_K$.  The checker in
\cref{sec:x3c-checker} operates on fixed-weight binary vectors and reads
$x_{\alpha_j}=\xi_j$ as the selection bits.  Hence every $\boldsymbol\xi$ must
admit a weight-$h$ extension $\mathbf x$ whose evaluation lies in one fixed
ideal coset.  This coset is independent of $\boldsymbol\xi$.  The ideal is
designed to have a large minimum, so that a nonzero ideal
element cannot produce an unintended short vector.  This section obtains such
an ideal from a lifted Reed--Solomon lattice~\cite{BennettPeikert23}, proves
that the required binary extensions exist even after their designated
coordinates are fixed, and establishes the additional norm bounds needed for
nonbinary representatives and scalar multiples.

\subsection{From lattice to ideal}

We first transfer an $\ell_2$-minimum bound for the lifted
Reed--Solomon lattice to a principal ideal of $\mathcal{O}_K$.

Fix an ordering $\F_q=\{a_1,\ldots,a_q\}$ and index the coordinates
of $\Z^q$ accordingly.  For $1\le k<q$, define
\[
 \mathbf H_q(k):=
 \begin{pmatrix}
  1&1&\cdots&1\\
  a_1&a_2&\cdots&a_q\\
  a_1^2&a_2^2&\cdots&a_q^2\\
  \vdots&\vdots&\ddots&\vdots\\
  a_1^{k-1}&a_2^{k-1}&\cdots&a_q^{k-1}
 \end{pmatrix}
 \in\F_q^{k\times q}.
\]
When $\mathbf H_q(k)$ is applied to
$\mathbf z\in\Z^q$, we first reduce the coordinates of $\mathbf z$ modulo
$q$.  The matrix $\mathbf H_q(k)$
defines the lifted Reed--Solomon lattice
\[
 \mathcal L_{q,k}:=\left\{\mathbf z\in\Z^q:
 \mathbf H_q(k)\mathbf z=\mathbf0\right\}.
\]

\begin{samepage}
\begin{lemma}[{\cite[Theorem~14]{BennettPeikert23}}]
\label{lem:rs-minimum}
Let $q$ be prime and $1\le k\le q/2$.  Then the first minimum of
$\mathcal L_{q,k}$ with respect to the $\ell_2$-norm $\|\cdot\|_2$ satisfies
\[
 \lambda_1(\mathcal L_{q,k})^2\ge 2k.
\]
\end{lemma}
\end{samepage}

Write
$\mathbf1_q^\perp:=\{\mathbf z\in\R^q:
\mathbf1_q^{\mathsf T}\mathbf z=0\}$.  To express the lattice equations and
evaluation at $\zeta$ in the same ring, associate $\mathbf z\in\Z^q$ with
\[
 f_{\mathbf z}:=\sum_{a\in\F_q}z_aX^a
 \in\Z[X]/(X^q-1).
\]
Every element of this quotient has a unique degree-$<q$ representative, which
we use to define its coefficients and its reduction $\bar f\in\F_q[X]$ modulo
$q$.  Let
\[
 \varphi:\Z[X]/(X^q-1)\longrightarrow \mathcal{O}_K,
 \qquad X\longmapsto\zeta,
\]
be the ring homomorphism given by evaluation at $\zeta$.  Its kernel is
generated by $\Phi_q(X)=1+X+\cdots+X^{q-1}$, and
$\varphi(f_{\mathbf z})=[\mathbf z]$.

\begin{theorem}
\label{thm:cyclotomic-descent}
For $1\le k\le q/2$,
\[
 \{[\mathbf z]:\mathbf z\in\mathcal L_{q,k}\}
 =\mathfrak a_{q,k}:=\pi^k\mathcal{O}_K.
\]
Moreover, the map $\mathbf z\mapsto[\mathbf z]$ restricts to an
additive-group isomorphism
\[
 \mathcal L_{q,k}\cap\mathbf1_q^\perp
 \ \longrightarrow\ \mathfrak a_{q,k}.
\]
Consequently, with respect to the canonical norm,
\[
 \lambda_1(\mathfrak a_{q,k})^2\ge2kq.
\]
\end{theorem}

\begin{proof}
 Since $k<q$, the polynomials
 $1,X,\ldots,X^{k-1}$ and
 $\binom X0,\binom X1,\ldots,\binom X{k-1}$ are related by an invertible
triangular change of basis over $\F_q$.  Hence the defining equations of
$\mathcal L_{q,k}$ are equivalent to
\[
 D^{[j]}\bar f_{\mathbf z}(1)=0
 \qquad(0\le j<k),
\]
because
$D^{[j]}\bar f_{\mathbf z}(1)=\sum_a z_a\binom aj$ in $\F_q$.
The multiplicity criterion in \cref{eq:hasse-multiplicity} and the identity
$X^q-1=(X-1)^q$ in $\F_q[X]$ therefore give
\[
 \{f_{\mathbf z}:\mathbf z\in\mathcal L_{q,k}\}
 =\bigl(q,(X-1)^k\bigr)
 \triangleleft\Z[X]/(X^q-1).
\]
Applying $\varphi$ gives the ideal $(q,\pi^k)$.  Total ramification gives
$(q)=(\pi^{q-1})$, so this image is
$(\pi^k)=\mathfrak a_{q,k}$ since $k\le q-1$.

Furthermore, $\Phi_q(X)\equiv(X-1)^{q-1}\pmod q$, so
$\ker(\varphi)=(\Phi_q)$ is contained in
$(q,(X-1)^k)$.  Since $\varphi((q,(X-1)^k))=\mathfrak a_{q,k}$, it follows that
\begin{equation}
\label{eq:ideal-preimage}
 \varphi^{-1}(\mathfrak a_{q,k})
 =\bigl(q,(X-1)^k\bigr)
 \quad\text{in }\Z[X]/(X^q-1).
\end{equation}

To prove surjectivity of the restricted map, fix
$v\in\mathfrak a_{q,k}$.  By \cref{eq:ideal-preimage}, choose
$f\in(q,(X-1)^k)$ with $\varphi(f)=v$.  Then $f(1)\in q\Z$, and
\[
 g:=f-\frac{f(1)}q\Phi_q
\]
remains in $(q,(X-1)^k)$, satisfies $\varphi(g)=v$, and has coefficient sum
zero.  Thus $g=f_{\mathbf z}$ for some
$\mathbf z\in\mathcal L_{q,k}\cap\mathbf1_q^\perp$.

For injectivity, suppose that
$\mathbf z\in\mathcal L_{q,k}\cap\mathbf1_q^\perp$ and $[\mathbf z]=0$.
The degree-$<q$ representative of $f_{\mathbf z}$ vanishes at $\zeta$, so it
equals $\ell\Phi_q$ for some $\ell\in\Z$.  Its coefficient sum is zero,
whereas $\Phi_q(1)=q$.  Thus $\ell=0$ and $\mathbf z=\mathbf0$.

Finally, every nonzero $v\in\mathfrak a_{q,k}$ has a unique nonzero
$\mathbf z\in\mathcal L_{q,k}\cap\mathbf1_q^\perp$ with $v=[\mathbf z]$.
By \cref{lem:cyclotomic-form,lem:rs-minimum},
\[
 \|v\|^2=q\|\mathbf z\|_2^2\ge2kq,
\]
which proves the minimum bound.
\end{proof}

\subsection{Binary representatives in ideal cosets}

For $\mathbf x\in\Z^q$, its \emph{syndrome}
$\mathbf H_q(k)\mathbf x\in\F_q^k$ has the coordinate indexed by
$j$ equal to
$\sum_{a\in\F_q}x_aa^j$ modulo $q$ for $0\le j<k$.  When $\mathbf x$ is
binary, these coordinates are the power sums of the positions $a$ for which
$x_a=1$.  \Cref{lem:coset-syndrome} identifies the vectors with fixed syndrome
and coefficient sum with one coset of $\mathfrak a_{q,k}$.

\begin{lemma}
\label{lem:coset-syndrome}
Assume $1\le k\le q/2$.  Let
$\mathbf u=(u_0,\ldots,u_{k-1})\in\F_q^k$, and let $s\in\Z$ have
residue $u_0$ modulo $q$.  There exists $\mathbf x_{\mathbf u,s}\in\Z^q$
such that
\[
 \mathbf H_q(k)\mathbf x_{\mathbf u,s}=\mathbf u,
 \qquad
 \mathbf1_q^{\mathsf T}\mathbf x_{\mathbf u,s}=s.
\]
Moreover, evaluation at $\zeta$ induces a bijection
\[
 \{\mathbf x\in\Z^q:
   \mathbf H_q(k)\mathbf x=\mathbf u,
   \ \mathbf1_q^{\mathsf T}\mathbf x=s\}
 \ \longrightarrow\
 [\mathbf x_{\mathbf u,s}]+\mathfrak a_{q,k},
 \qquad \mathbf x\longmapsto[\mathbf x].
\]
Thus every element of this ideal coset has a unique representative with
syndrome $\mathbf u$ and coefficient sum $s$.
\end{lemma}

\begin{proof}
A Vandermonde $k\times k$ minor of $\mathbf H_q(k)$ has determinant
$\prod_{1\le i<j\le k}(a_j-a_i)\ne0$ in $\F_q$.  Hence the matrix has
rank $k$.  Choose $\mathbf x_0\in\Z^q$ whose reduction modulo $q$
satisfies $\mathbf H_q(k)\mathbf x_0=\mathbf u$ over $\F_q$.
Its coefficient sum is congruent to $u_0$, and hence to $s$ modulo $q$.
The vector $\mathbf1_q$ belongs to $\mathcal L_{q,k}$: its zeroth power sum is
$q=0$ in $\F_q$, and
$\sum_{a\in\F_q}a^j=0$ for $1\le j<k<q-1$.  Moreover,
$[\mathbf1_q]=0$.  Therefore adding a suitable integer multiple of
$\mathbf1_q$ to $\mathbf x_0$ gives a vector $\mathbf x_{\mathbf u,s}$ with
the required coefficient sum without changing its syndrome or its image at
$\zeta$.

Every other vector with syndrome $\mathbf u$ and coefficient sum $s$ differs
from $\mathbf x_{\mathbf u,s}$ by an element of
$\mathcal L_{q,k}\cap\mathbf1_q^\perp$.  The restricted isomorphism in
\cref{thm:cyclotomic-descent} now gives the stated bijection.
\end{proof}

We use the following parameters to specialize
\cref{prop:wan-power-sum-estimates}.
For every sufficiently large prime $q$, put
\begin{equation}
\label{eq:rs-parameters}
 k:=\left\lfloor\frac{q^{1/1000}}2\right\rfloor,
 \qquad
 h:=\left\lfloor\frac{201k}{200}\right\rfloor,
 \qquad
 T:=\lfloor q^{1/2000}\rfloor.
\end{equation}
These integers can be computed without real-number
approximations.  The value $k$ is the largest nonnegative integer satisfying
$(2k)^{1000}\le q$, and $T$ is the largest nonnegative integer satisfying
$T^{2000}\le q$.  Fixed-degree integer-root algorithms compute $k$ and $T$
in time polynomial in $\log q$, after which integer division gives $h$.
Here $k$ is the number of rows of $\mathbf H_q(k)$.  For each
vector $\mathbf x\in\{0,1\}^{\F_q}$ considered below, $h$ is its weight, and $T$
is the maximum number of its coordinates whose values may be fixed in
advance.  The relevant scale relations are
$T=o(k)$, $h=(201/200)k+O(1)$, and $h=o(q)$.  Set
\begin{equation}
\label{eq:HL-parameters}
 H:=h(q-h),
 \qquad
 L:=2kq.
\end{equation}
By \cref{lem:cyclotomic-form}, every
$\mathbf x\in\{0,1\}^{\F_q}$ of weight $h$ satisfies $\|[\mathbf x]\|^2=H$.
By \cref{thm:cyclotomic-descent}, every nonzero element of
$\mathfrak a_{q,k}$ has canonical squared norm at least $L$.

For a prescribed syndrome $\mathbf u=(u_0,\ldots,u_{k-1})$, let
$P\subseteq\F_q$ be the \emph{position set} indexing the coordinates fixed
in advance, and let $\boldsymbol\eta\in\{0,1\}^P$ specify their values.
Define
\[
 \mathcal F(\mathbf u;P,\boldsymbol\eta):=
 \{\mathbf x\in\{0,1\}^{\F_q}:\mathbf H_q(k)\mathbf x=\mathbf u,\
 \mathbf1_q^{\mathsf T}\mathbf x=h,\
 \mathbf x|_P=\boldsymbol\eta\}.
 \]
The first row of $\mathbf H_q(k)$ forces $u_0=h$ in $\F_q$.

\begin{theorem}
\label{thm:uniform-completions}
For every sufficiently large prime $q$, every $\mathbf u\in\F_q^k$ with
$u_0=h$, every position set $P\subseteq\F_q$ with $|P|\le T$, and every
$\boldsymbol\eta\in\{0,1\}^P$, put $t:=\sum_{a\in P}\eta_a$.  Then
\begin{equation}
\label{eq:uniform-completion-bounds}
 |\mathcal F(\mathbf u;P,\boldsymbol\eta)|
 =\Theta\left(\frac{q^{h-t-k+1}}{(h-t)!}\right).
\end{equation}
 Both the threshold for $q$ and the implied constants are uniform in
 $\mathbf u$, $P$, and $\boldsymbol\eta$.
\end{theorem}

\begin{proof}
Fix $\mathbf u$, $P$, and $\boldsymbol\eta$, and put $m_0:=h-t$.  By
\cref{eq:rs-parameters}, all sufficiently large $q$ satisfy
\[
 k+2\le m_0<q-2.
\]
The fixed coordinates contain exactly $t$ ones.  Subtract their contribution
by defining
\begin{equation}
\label{eq:shifted-syndrome}
 u'_j:=u_j-\sum_{a\in P}\eta_aa^j
 \qquad(1\le j<k).
\end{equation}
Every $\mathbf x\in\mathcal F(\mathbf u;P,\boldsymbol\eta)$ has exactly
$m_0$ coordinates outside $P$ equal to one.  Ordering their positions gives
a tuple $(z_1,\ldots,z_{m_0})\in(\F_q\setminus P)^{m_0}$ such that
$z_i\ne z_\ell$ for every $i\ne\ell$ and
\begin{equation}
\label{eq:ordered-power-sums}
 \sum_{i=1}^{m_0}z_i^j=u'_j
 \qquad(1\le j<k).
\end{equation}
Conversely, every pairwise distinct tuple in
$(\F_q\setminus P)^{m_0}$ satisfying \cref{eq:ordered-power-sums}
determines a unique vector by setting $x_{z_i}=1$, setting
$\mathbf x|_P=\boldsymbol\eta$, and setting all remaining coordinates to
zero.  The equation for $j=0$ follows from $u_0=h$.  Hence the number of
these ordered tuples is $m_0!\,|\mathcal F(\mathbf u;P,\boldsymbol\eta)|$.

Let $A$ count the tuples satisfying
\cref{eq:ordered-power-sums} without requiring their entries to be
pairwise distinct or to lie outside $P$.  After the shift in
\cref{eq:shifted-syndrome}, part~(i) of
\cref{prop:wan-power-sum-estimates} gives
$A=q^{m_0-k+1}(1+o(1))$.  We bound the excluded tuples in two cases.
\par\noindent
\emph{Case I (collision).}  For fixed $i<\ell$,
part~(ii) of \cref{prop:wan-power-sum-estimates} shows that
the number of solutions with
$z_i=z_\ell$ is $O(q^{m_0-k})$.
\par\noindent
\emph{Case II (entry in $P$).}  For fixed $i$ and $a\in P$, substituting
$z_i=a$ in \cref{eq:ordered-power-sums} gives
$\sum_{r\ne i}z_r^j=u'_j-a^j$ for $1\le j<k$.
Part~(iii) of \cref{prop:wan-power-sum-estimates} shows that
the number of solutions with $z_i=a$ is
$O(q^{m_0-k})$.

The explicit error bounds in
\cref{prop:wan-power-sum-estimates} are uniform in the prescribed
power-sum vector $(u'_j)_{1\le j<k}$.  For the
parameters in \cref{eq:rs-parameters}, each relative error is at
most
\[
 (2k)^{m_0}q^{-(m_0-k-2)/2}=q^{-\Omega(k)}=o(1).
\]
Thus the estimate for $A$, the bound for $z_i=z_\ell$, and the
bound for $z_i=a$ are uniform over all shifted power-sum vectors
$(u'_j)_{1\le j<k}$ and all
choices of $i$, $\ell$, and $a$.

The union bound over $\binom{m_0}{2}$ equalities $z_i=z_\ell$ and
$m_0|P|$ equalities $z_i=a$ with $a\in P$ shows that the nonnegative
difference satisfies
\[
 0\le A-m_0!\,|\mathcal F(\mathbf u;P,\boldsymbol\eta)|
 =O\bigl((m_0^2+m_0|P|)q^{m_0-k}\bigr).
\]
After division by the main term $q^{m_0-k+1}$, this error is
$O((m_0^2+m_0|P|)/q)=o(1)$ by the scale relations following
\cref{eq:rs-parameters}.  Together with the estimate for $A$, this gives
\[
 m_0!\,|\mathcal F(\mathbf u;P,\boldsymbol\eta)|
 =q^{m_0-k+1}(1+o(1)).
\]
Substituting $m_0=h-t$ proves \cref{eq:uniform-completion-bounds}.
\end{proof}

\begin{corollary}
\label{cor:fixed-coordinate-ratio}
Let $\mathbf u,P,\boldsymbol\eta$, and $t$ be as in
\cref{thm:uniform-completions}.  For $r\in\{1,2\}$, let
$Q\subseteq\F_q\setminus P$ have size $r$ and satisfy $|P|+r\le T$.
Put $P':=P\cup Q$ and extend $\boldsymbol\eta$ to
$\boldsymbol\eta'\in\{0,1\}^{P'}$ by setting $\eta'_a=1$ for every $a\in Q$.
Then
\begin{equation}
\label{eq:fixed-coordinate-ratio}
 \frac{|\mathcal F(\mathbf u;P',\boldsymbol\eta')|}
 {|\mathcal F(\mathbf u;P,\boldsymbol\eta)|}
 =O\left((h/q)^r\right).
\end{equation}
The implied constant is uniform in $\mathbf u$, $P$, $\boldsymbol\eta$, and
$Q$.
\end{corollary}

\begin{proof}
For all sufficiently large $q$,
$t+r\le|P|+r\le T<h$, so $h-t-r\ge0$.  The numerator in
\cref{eq:fixed-coordinate-ratio} counts exactly the vectors in
$\mathcal F(\mathbf u;P,\boldsymbol\eta)$ that equal one at every coordinate
in $Q$.  Applying \cref{thm:uniform-completions} to the numerator and
denominator gives
\[
 \frac{|\mathcal F(\mathbf u;P',\boldsymbol\eta')|}
 {|\mathcal F(\mathbf u;P,\boldsymbol\eta)|}
 =O\left(q^{-r}\frac{(h-t)!}{(h-t-r)!}\right)
 =O\left((h/q)^r\right).
\]
\end{proof}

\subsection{Norm separation}

The preceding count supplies the intended binary representatives.  To use
their ideal coset in the module construction, we must rule out shorter
representatives with nonbinary coefficients and short points arising from
nonzero scalar multiples.

\begin{lemma}
\label{lem:binary-shell}
Let $\mathbf u\in\F_q^k$ satisfy $u_0=h$, and let
$\mathbf x'\in\Z^q$ have syndrome $\mathbf u$ and coefficient sum $h$.
Put $c':=[\mathbf x']$.  Every $v\in c'+\mathfrak a_{q,k}$ has a unique
integral coefficient representative $\mathbf z\in\Z^q$ satisfying
\[
 [\mathbf z]=v,
 \qquad
 \mathbf H_q(k)\mathbf z=\mathbf u,
 \qquad
 \sum_a z_a=h.
\]
If $\mathbf z\in\{0,1\}^{\F_q}$, then it has weight $h$ and
$\|v\|^2=H$.  Otherwise, $\|v\|^2\ge H+2q$.
\end{lemma}

\begin{proof}
The existence and uniqueness of $\mathbf z$ follow directly from
\cref{lem:coset-syndrome}.  By \cref{lem:cyclotomic-form},
\[
 \|v\|^2=q\sum_a z_a^2-h^2
 =H+q\sum_a z_a(z_a-1).
\]
For every integer $z_a$, the quantity $z_a(z_a-1)$ is a nonnegative even
integer.  The sum is zero exactly when every $z_a\in\{0,1\}$, and is otherwise
at least two.  This proves the stated claim.
\end{proof}

We next choose $c$ so that $wc$ remains separated from
$\mathfrak a_{q,k}$ for every nonzero $w\in\mathcal{O}_K$ with
$\|w\|^2\le3q$.  Define $c\in\mathcal{O}_K$ and
$c+\mathfrak a_{q,k}$ by
\begin{equation}
\label{eq:robust-center}
 m:=\left\lfloor\frac k2\right\rfloor,
 \qquad
 c:=h+\pi^m.
\end{equation}
The scalar summand $h$ fixes the coefficient sum at $h$, since
$h+(1-X)^m$ evaluates to $h$ at $X=1$.  The term $\pi^m$, with
$m=\lfloor k/2\rfloor$, balances the two orders of vanishing, $m$ and $k-m$,
that arise from the $\pi^k$-divisibility condition.  Thus each of the two
sparse-polynomial arguments below retains $k/2+O(1)$ orders of vanishing.
Let $\mathbf c\in\Z^q$ be the coefficient vector of $h+(1-X)^m$.  Then
$[\mathbf c]=c$, and its syndrome
\[
 \mathbf u_c:=\mathbf H_q(k)\mathbf c\in\F_q^k
\]
satisfies $(u_c)_0=h$.  Thus \cref{lem:coset-syndrome} identifies the vectors
of syndrome $\mathbf u_c$ and coefficient sum $h$ with the coset
$c+\mathfrak a_{q,k}$.

The proof of \cref{thm:robust-center} combines a sparse integral
coefficient representative with a bound on the multiplicity of a sparse
polynomial at $X=1$.  The next lemma records both facts.
\begin{lemma}
\label{lem:sparse-tools}
The following statements hold.
\begin{enumerate}[label=(\roman*)]
 \item If $v\in\mathcal{O}_K$ and $\|v\|^2<q^2/4$, then $v$ has
 an integral coefficient representative with at most
 $2\|v\|^2/q$ nonzero coefficients.
 \item If
 \[
  F(X)=\sum_{i=1}^s c_iX^{a_i}\in\F_q[X]
 \]
 is nonzero, where every $c_i\ne0$ and the exponents
 $a_i\in\{0,\ldots,q-1\}$ are distinct, then its order of
 vanishing at
 $X=1$ is at most $s-1$.
\end{enumerate}
\end{lemma}

\begin{proof}
For (i), choose a coefficient vector
$\mathbf z=(z_a)_{a\in\F_q}\in\Z^q$ with $[\mathbf z]=v$.  If no value occurs
among the coordinates $z_a$ more than $q/2$ times, then at least $q^2/4$
unordered pairs have unequal coordinates, contradicting the last expression
in \cref{lem:cyclotomic-form}.  Subtract the most frequent coordinate value
from every $z_a$.  This does not change $[\mathbf z]=v$.  Let $s$ be the
number of nonzero coefficients in the resulting representative.  The
preceding argument gives $s<q/2$.  Each pair formed
by a nonzero coordinate and a zero coordinate gives
\[
 \|v\|^2\ge s(q-s)\ge\frac{sq}{2},
\]
which proves (i).

For (ii), multiplicity at least $s$ would make the first $s$ Hasse derivatives
vanish at one.  The resulting homogeneous system has determinant
\[
 \det\left(\binom{a_i}{j}\right)_{
       0\le j<s,\ 1\le i\le s}
 =\frac{\prod_{1\le i<\ell\le s}(a_\ell-a_i)}
        {\prod_{j=0}^{s-1}j!},
\]
which is nonzero in $\F_q$: the exponents are distinct, and $s\le q$ makes
every factorial in the denominator nonzero.  Hence every $c_i$ would vanish,
a contradiction.
\end{proof}

With $b:=1/100$, the next theorem separates every short nonzero
multiple of $c$ from $\mathfrak a_{q,k}$.
\begin{theorem}
\label{thm:robust-center}
For all sufficiently large primes $q$, every nonzero
$w\in\mathcal{O}_K$ with $\|w\|^2\le3q$ satisfies
\[
 \dist(wc,\mathfrak a_{q,k})^2>bL.
\]
\end{theorem}

\begin{proof}
We prove the theorem by contradiction.  Suppose that there exist a nonzero
$w\in\mathcal{O}_K$ with $\|w\|^2\le3q$ and an element
$e\in wc+\mathfrak a_{q,k}$ such that
$\|e\|^2\le bL=2bkq$.  For sufficiently large $q$, both $2bkq$ and $3q$ are
below $q^2/4$.  By part~(i) of \cref{lem:sparse-tools}, choose integral
coefficient representatives $\mathbf z_e$ and $\mathbf z_w$ of $e$ and $w$
with at most $4bk$ and $6$ nonzero coefficients, respectively.  Put
$f_e:=f_{\mathbf z_e}$ and $f_w:=f_{\mathbf z_w}$.  Since
$\varphi(h+(1-X)^m)=c$ and $e-wc\in\mathfrak a_{q,k}$,
\cref{eq:ideal-preimage} gives
\[
 f_e-f_w\bigl(h+(1-X)^m\bigr)\in\bigl(q,(X-1)^k\bigr).
\]
Reducing modulo $q$ and using $X^q-1=(X-1)^q$ in $\F_q[X]$ gives
\begin{equation}
\label{eq:center-congruence}
 \bar f_e-h\bar f_w
 \equiv\bar f_w(1-X)^m\pmod{(X-1)^k}.
\end{equation}
In particular, $\bar f_e-h\bar f_w$ is divisible by $(X-1)^m$.
For sufficiently large $k$, this polynomial has at most $4bk+6\le m$ nonzero
coefficients.  Part~(ii) of \cref{lem:sparse-tools} therefore forces
$\bar f_e-h\bar f_w=0$ in $\F_q[X]$.

It follows from \cref{eq:center-congruence} that
$\bar f_w$ has a zero of order at least $k-m$ at $X=1$.  This polynomial is
nonzero.  Otherwise, $\mathbf z_w=q\mathbf z'$ for some $\mathbf z'\in\Z^q$.
Since $[\mathbf z']=w/q\ne0$, the vector $\mathbf z'$ is not constant.  At
least $q-1$ pairs of its coordinates therefore differ, and
\cref{lem:cyclotomic-form} gives
\[
 \|w\|^2=q^2\|[\mathbf z']\|^2\ge q^2(q-1)>3q.
\]
Since $\bar f_w$ has at most six nonzero coefficients, part~(ii) of
\cref{lem:sparse-tools} bounds its order of vanishing at $X=1$ by $5$.  This
is impossible because $k-m\to\infty$.
\end{proof}

We finally record the norm gap between roots of unity and all other
nonzero elements of $\mathcal{O}_K$.
\begin{lemma}
\label{lem:second-shell}
For every nonzero $w\in\mathcal{O}_K$, either
$w=\pm\zeta^j$ for some $j$, in which case $\|w\|^2=q-1$, or~
$\|w\|^2\ge2q-4$.
\end{lemma}

\begin{proof}
If $\|w\|^2\ge q^2/4$, then $\|w\|^2\ge2q-4$.  Otherwise,
part~(i) of
\cref{lem:sparse-tools} gives an integral coefficient representative with
exactly $s<q/2$ nonzero coefficients.  If $s\ge2$, every pair formed by one
nonzero coordinate and one zero coordinate contributes to
\[
 \|w\|^2\ge s(q-s)\ge2(q-2).
\]
If $s=1$, then $w=t\zeta^j$ for a nonzero integer $t$.  The cases
$t=\pm1$ are precisely the roots of unity, whereas $|t|\ge2$ gives squared
norm at least $4(q-1)$.
\end{proof}
